%% This is file `elsarticle-template-1-num.tex',
%%
%% Copyright 2009 Elsevier Ltd
%%
%% This file is part of the 'Elsarticle Bundle'.
%% ---------------------------------------------

\documentclass[a4paper]{article}

\usepackage{graphicx}
\usepackage{import}
\usepackage{xifthen}
\usepackage{pdfpages}
\usepackage{transparent}
\usepackage{cancel}
\usepackage{amssymb}
\usepackage{amsmath}
\usepackage{amsthm}
\usepackage{ mathrsfs }
\usepackage{ dsfont }
\usepackage{tikz-cd}
\usepackage[a4paper,top=3cm,bottom=2cm,left=3cm,right=3cm,marginparwidth=1.75cm]{geometry}
\usepackage{float}
\usepackage{lineno}
\usepackage{hyperref}
\usepackage{enumitem}

\newtheorem{name}{Printed output}
\newtheorem{mydef}{Definition}
\newtheoremstyle{break}%
  {}%
  {}%
  {}%
  {}%
  {\bfseries}%
  {.}%
  {\newline}%
  {}

\theoremstyle{plain}
\newtheorem{question}{Question}

\newtheorem{theorem}{Theorem}[section]
\newtheorem*{lemma*}{Lemma}
\newtheorem{lemma}[theorem]{Lemma}
\newtheorem{corollary}[theorem]{Corollary}
\newtheorem{proposition}[theorem]{Proposition}

\theoremstyle{definition}
\newtheorem{definition}[theorem]{Definition}
\newtheorem*{note*}{Note}
\newtheorem*{claim}{Claim}
\newtheorem*{remark}{Remark}
\newtheorem{example}{Example}
\newtheorem*{solution}{Solution}
\newtheorem{exercise}{Exercise}

\numberwithin{equation}{subsection}

\newcommand{\incfig}[2][1]{%
    \def\svgwidth{#1\columnwidth}%
    \import{./Figures/}{#2.pdf_tex}%
}
\newcommand{\hodge}{{*}}
\newcommand{\tang}{\operatorname{\textbf{\textup{t}}}}
\newcommand{\norm}{\operatorname{\textbf{\textup{n}}}}
\newcommand{\R}{\mathbb{R}}
\newcommand{\bP}{\mathbb{bP}}
\newcommand{\N}{\mathbb{N}}
\newcommand{\Q}{\mathbb{Q}}
\newcommand{\A}{\mathbb{A}}
\newcommand{\bS}{\mathbb{S}}
\newcommand{\Z}{\mathbb{Z}}
\newcommand{\F}{\mathbb{F}}
\newcommand{\B}{\mathbb{B}}
\newcommand{\T}{\mathbb{T}}
\newcommand{\K}{\mathbb{K}}
\newcommand{\C}{\mathbb{C}}
\newcommand{\D}{\mathbb{D}}

\usepackage{comment}
\newcommand{\bmv}{\bm{v}}
\newcommand{\bmu}{\bm{u}}
\newcommand{\bmw}{\bm{w}}
\newcommand{\bmx}{\bm{x}}

\newcommand{\cF}{\mathcal{F}}
\newcommand{\cA}{\mathcal{A}}
\newcommand{\cM}{\mathcal{M}}
\newcommand{\cU}{\mathcal{U}}
\newcommand{\cB}{\mathcal{B}}
\newcommand{\cD}{\mathcal{D}}
\newcommand{\cL}{\mathcal{L}}
\newcommand{\cS}{\mathcal{S}}
\newcommand{\cG}{\mathcal{G}}
\newcommand{\cT}{\mathcal{T}}
\newcommand{\cC}{\mathcal{C}}
\newcommand{\cH}{\mathcal{H}}
\newcommand{\cO}{\mathcal{O}}
\newcommand{\cJ}{\mathcal{J}}

\newcommand{\frp}{\mathfrak{p}}
\newcommand{\frq}{\mathfrak{q}}
\newcommand{\frm}{\mathfrak{m}}
\newcommand{\frP}{\mathfrak{P}}

\newcommand{\msF}{\mathscr{F}}
\newcommand{\msA}{\mathscr{A}}
\newcommand{\msM}{\mathscr{M}}
\newcommand{\msE}{\mathscr{E}}
\newcommand{\msU}{\mathscr{U}}
\newcommand{\msB}{\mathscr{B}}
\newcommand{\msD}{\mathscr{D}}
\newcommand{\msL}{\mathscr{L}}
\newcommand{\msS}{\mathscr{S}}
\newcommand{\msG}{\mathscr{G}}
\newcommand{\msT}{\mathscr{T}}
\newcommand{\msC}{\mathscr{C}}
\newcommand{\msH}{\mathscr{H}}
\newcommand{\msO}{\mathscr{O}}
\newcommand{\msJ}{\mathscr{J}}
\newcommand{\msP}{\mathscr{P}}

\newcommand{\vect}[1]{\underline{\textbf{#1}}}
\newcommand{\llangle}{\left\langle}
\newcommand{\rrangle}{\right\rangle}
\newcommand{\llbrace}{\left\lbrace}
\newcommand{\rrbrace}{\right\rbrace}
\newcommand{\setof}[1]{\left\lbrace #1 \right\rbrace}

\newcommand{\cFA}{\cF(\cA)}
\newcommand{\sigmaA}{\sigma(\cA)}
\newcommand{\del}[2]{\frac{\partial #1}{\partial #2}}

\newcommand{\innerproductb}[2]{\llangle\bm{#1},\bm{#2}\rrangle}
\newcommand{\innerproduct}[2]{\llangle{#1},{#2}\rrangle}

\newcommand{\cl}{\operatorname{cl}}
\newcommand{\coker}{\operatorname{coker}}
\newcommand{\re}{\operatorname{Re}}
\newcommand{\im}{\operatorname{Im}}
\newcommand{\dist}{\operatorname{dist}}
\newcommand{\Span}{\operatorname{Span}}
\newcommand{\Supp}{\operatorname{Supp}}
\newcommand{\id}{\operatorname{id}}
\newcommand{\sh}{\operatorname{sh}}
\newcommand{\var}{\operatorname{Var}}
\newcommand{\Proj}{\operatorname{Proj}}
\newcommand{\Ker}{\operatorname{Ker}}
\newcommand{\Mor}{\operatorname{Mor}}
\newcommand{\supp}{\operatorname{supp}}
\newcommand{\Ran}{\operatorname{Ran}}
\newcommand{\Res}{\operatorname{Res}}
\newcommand{\Mod}{\operatorname{Mod}}
\newcommand{\res}{\operatorname{res}}
\newcommand{\Frac}{\operatorname{Frac}}
\newcommand{\End}{\operatorname{End}}
\newcommand{\Sym}{\operatorname{Sym}}
\newcommand{\Spec}{\operatorname{Spec}}
\newcommand*{\sheafhom}{\mathcal{H}\kern -.5pt \it{ om}}
\newcommand*{\sheafspec}{\mathcal{S}\kern -.5pt \it{ pec}}
\newcommand*{\sheafproj}{{\mathcal{P}\kern -.5pt \it{ roj}}}
\newcommand{\Div}{\operatorname{div}}
\newcommand{\Aut}{\operatorname{Aut}}
\newcommand{\Hom}{\operatorname{Hom}}
\newcommand{\Qcoh}{\operatorname{Qcoh}}
\newcommand{\PGL}{\operatorname{PGL}}
\newcommand{\Op}{\operatorname{Op}}
\newcommand{\WF}{\operatorname{WF}}
\newcommand{\Char}{\operatorname{Char}}
\newcommand{\Ell}{\operatorname{Ell}}
\newcommand{\RE}{\operatorname{Re}}
\newcommand{\IM}{\operatorname{Im}}
\newcommand{\cosec}{\operatorname{cosec}}

\setlength{\parindent}{0pt}
\setlength{\parskip}{1em}

\begin{document}

\title{\textbf{Problem Sheet: Trig Substitutions 1 Solution}}
\author{\textbf{Jeffrey Thompson}}
\maketitle

\begin{exercise}
Using a trigonometric substitution, find the indefinite integral
\[
    \int \sqrt{1 - 4x^2} \, dx
\]
Then state the interval of $x$ values where this substitution is valid.
\end{exercise}

\begin{solution}

The identities that look like $1-4x^2$ are
\[
    1-\sin^2\theta=\cos^2\theta, \qquad 1-\cos^2\theta=\sin^2\theta
\]
so we use a scaled version of the first and let
\[
    2x = \sin \theta
\]
with
\[
    -\frac{\pi}{2} \leq \theta \leq \frac{\pi}{2}.
\]
This is a convenient choice because then $\cos\theta \geq 0$.

Therefore
\[
    x = \frac{1}{2}\sin \theta
\]
and
\[
    \frac{dx}{d\theta} = \frac{1}{2}\cos \theta
\]
so
\[
    dx = \frac{1}{2}\cos \theta \, d\theta.
\]

Now substitute from left to right:
\begin{align*}
    \int \sqrt{1 - 4x^2} \, dx
    &= \int \sqrt{1 - \sin^2 \theta} \, dx \tag{using $2x = \sin \theta$}\\
    &= \int \sqrt{\cos^2 \theta}\cdot \frac{1}{2}\cos \theta \, d\theta \tag{using $dx = \frac{1}{2}\cos \theta\,d\theta$}\\
    &= \frac{1}{2}\int \cos^2 \theta \, d\theta
\end{align*}

Here we used that
\[
    \sqrt{\cos^2\theta} = \cos\theta,
\]
which is valid on our chosen interval because $\cos\theta \geq 0$ there.

Now use
\[
    \cos^2 \theta = \frac{1+\cos 2\theta}{2}
\]
to find
\begin{align*}
    \frac{1}{2}\int \cos^2 \theta \, d\theta
    &= \frac{1}{2}\int \frac{1+\cos 2\theta}{2}\,d\theta\\
    &= \frac{1}{4}\int 1\,d\theta + \frac{1}{4}\int \cos 2\theta\,d\theta\\
    &= \frac{1}{4}\theta + \frac{1}{8}\sin 2\theta + c
\end{align*}

Now return to $x$. Since $2x = \sin \theta$, we have
\[
    \theta = \sin^{-1}(2x)
\]
and also
\[
    \sin 2\theta = 2\sin \theta \cos \theta.
\]
Therefore
\begin{align*}
    \frac{1}{4}\theta + \frac{1}{8}\sin 2\theta + c
    &= \frac{1}{4}\theta + \frac{1}{8}\left(2\sin \theta \cos \theta\right) + c\\
    &= \frac{1}{4}\theta + \frac{1}{4}\sin \theta \cos \theta + c\\
    &= \frac{1}{4}\theta + \frac{1}{4}(2x)\sqrt{1-\sin^2\theta} + c\\
    &= \frac{1}{4}\theta + \frac{1}{2}x\sqrt{1-4x^2} + c\\
    &= \frac{1}{4}\sin^{-1}(2x) + \frac{1}{2}x\sqrt{1-4x^2} + c
\end{align*}

Hence
\[
    \int \sqrt{1 - 4x^2} \, dx
    = \frac{1}{2}x\sqrt{1-4x^2} + \frac{1}{4}\sin^{-1}(2x) + c
\]

When using this substitution, we are working in the range
\[
    -1 \leq 2x \leq 1
\]
which gives
\[
    -\frac{1}{2} \leq x \leq \frac{1}{2}.
\]
This is also exactly the range where the square root term $\sqrt{1-4x^2}$ is defined.

\qed
\end{solution}

% ------------------------------------------------------------

\begin{exercise}
Using a trigonometric substitution, compute
\[
    \int_{\frac{1}{2}}^{1}\frac{1}{x\sqrt{1-x^2}}\,dx.
\]
\end{exercise}

\begin{solution}

Since the integrand contains $\sqrt{1-x^2}$, we use
\[
    x = \cos \theta.
\]
Since the bounds satisfy $\frac{1}{2}\leq x \leq 1$, we take
\[
    0 \leq \theta \leq \frac{\pi}{2},
\]
so that both $\cos\theta \geq 0$ and $\sin\theta \geq 0$.

Then
\begin{align*}
    x &= \cos \theta\\
    dx &= -\sin \theta\,d\theta\\
    \sqrt{1-x^2} &= \sqrt{1-\cos^2\theta} = \sin \theta
\end{align*}
where the final step is valid because $\sin \theta \geq 0$ on the interval we use.

Now convert the bounds:
\begin{align*}
    x = \frac{1}{2} &\therefore \cos \theta = \frac{1}{2} \therefore \theta = \frac{\pi}{3}\\
    x = 1 &\therefore \cos \theta = 1 \therefore \theta = 0
\end{align*}

Substitute from left to right:
\begin{align*}
    \int_{\frac{1}{2}}^{1}\frac{1}{x\sqrt{1-x^2}}\,dx
    &= \int_{x=\frac{1}{2}}^{x=1}\frac{1}{x\sqrt{1-x^2}}\,dx\\
    &= \int_{\theta=\frac{\pi}{3}}^{\theta=0}\frac{1}{\cos\theta\sqrt{1-\cos^2\theta}}(-\sin\theta)\,d\theta\\
    &= \int_{\theta=\frac{\pi}{3}}^{\theta=0}\frac{-\sin\theta}{\cos\theta\sin\theta}\,d\theta\\
    &= \int_{\theta=\frac{\pi}{3}}^{\theta=0}-\sec\theta\,d\theta\\
    &= \int_{0}^{\frac{\pi}{3}}\sec\theta\,d\theta
\end{align*}

Now use the formula book:
\[
    \int \sec\theta\,d\theta = \ln|\sec\theta + \tan\theta| + c
\]
Therefore
\begin{align*}
    \int_{0}^{\frac{\pi}{3}}\sec\theta\,d\theta
    &= \left[\ln|\sec\theta + \tan\theta|\right]_{\theta=0}^{\theta=\frac{\pi}{3}}\\
    &= \ln\left(\sec\frac{\pi}{3} + \tan\frac{\pi}{3}\right) - \ln(\sec 0 + \tan 0)
\end{align*}

Evaluate the trig values:
\[
    \sec\frac{\pi}{3} = 2, \qquad \tan\frac{\pi}{3} = \sqrt{3}, \qquad \sec 0 = 1, \qquad \tan 0 = 0.
\]
So
\begin{align*}
    \int_{\frac{1}{2}}^{1}\frac{1}{x\sqrt{1-x^2}}\,dx
    &= \ln(2+\sqrt{3}) - \ln(1)\\
    &= \ln(2+\sqrt{3})
\end{align*}

\qed
\end{solution}

% ------------------------------------------------------------

\begin{exercise}
Compute the following indefinite integral, assuming that $4-x^2\geq 0$ and $x\neq0$.

\[
    \int \frac{1}{x^2 \sqrt{4 - x^2}}dx
\]
\end{exercise}

\begin{solution}

We look for a trig identity that matches $4-x^2$. Since
\begin{align*}
    \sin^2\theta + \cos^2\theta &= 1\\
    4\sin^2\theta + 4\cos^2\theta &= 4\\
    4\sin^2\theta &= 4 - 4\cos^2\theta
\end{align*}
we choose
\[
    x = 2\cos\theta.
\]
We take
\[
    0 \leq \theta \leq \pi
\]
so that $\sin\theta \geq 0$, which will make the square root step below valid.

Then
\begin{align*}
    x^2 &= 4\cos^2\theta\\
    \sqrt{4-x^2} &= \sqrt{4-4\cos^2\theta} = 2\sin\theta\\
    \frac{dx}{d\theta} &= -2\sin\theta\\
    dx &= -2\sin\theta\,d\theta
\end{align*}

Now substitute from left to right:
\begin{align*}
    \int \frac{1}{x^2 \sqrt{4 - x^2}}dx
    &= \int \frac{1}{4\cos^2\theta\sqrt{4-x^2}}dx \tag{using $x^2 = 4\cos^2\theta$}\\
    &= \int \frac{1}{4\cos^2\theta \cdot 2\sin\theta}dx \tag{using $\sqrt{4-x^2} = 2\sin\theta$}\\
    &= \int \frac{1}{4\cos^2\theta \cdot \cancel{2\sin\theta}}(-\cancel{2\sin\theta})\,d\theta \tag{using $dx = -2\sin\theta\,d\theta$}\\
    &= -\frac{1}{4}\int \sec^2\theta\,d\theta
\end{align*}

Using the formula book,
\[
    \int \sec^2\theta\,d\theta = \tan\theta + c
\]
so
\[
    \int \frac{1}{x^2 \sqrt{4 - x^2}}dx = -\frac{1}{4}\tan\theta + c.
\]

Now write everything back in terms of $x$. Since $x = 2\cos\theta$, we have
\[
    \cos\theta = \frac{x}{2}
\]
and therefore
\begin{align*}
    \tan\theta
    &= \frac{\sin\theta}{\cos\theta}\\
    &= \frac{\sqrt{1-\cos^2\theta}}{\cos\theta}\\
    &= \frac{\sqrt{1-\left(\frac{x}{2}\right)^2}}{\frac{x}{2}}\\
    &= \frac{\sqrt{\frac{4-x^2}{4}}}{\frac{x}{2}}\\
    &= \frac{\frac{1}{2}\sqrt{4-x^2}}{\frac{x}{2}}\\
    &= \frac{\sqrt{4-x^2}}{x}
\end{align*}

Hence
\[
    \int \frac{1}{x^2 \sqrt{4 - x^2}}dx
    = -\frac{1}{4}\cdot \frac{\sqrt{4-x^2}}{x} + c
\]
that is,
\[
    \int \frac{1}{x^2 \sqrt{4 - x^2}}dx
    = -\frac{\sqrt{4-x^2}}{4x} + c.
\]

Since the original integrand contains $x^2$ in the denominator, we also require
\[
    x \neq 0.
\]

\qed
\end{solution}

\end{document}
